% =============================================================
%  Chapter — HVDC Topologies
%  Assets (SVG source + PDF build) live in this same subfolder.
% =============================================================
\sectionslide{AC Loadflow}{AC loadflow calculations}


% ---- Symbols ----
\begin{frame}{Symbols}
  \renewcommand{\arraystretch}{1.3}
  \begin{tabular}{@{}p{1.3cm}p{9.0cm}@{}}
    \textbf{Sym.} & \textbf{Description} \\
    \hline
    $P$        & Active power \\
    $V_{LL}$   & Line-to-line (phase-phase) voltage \\
    $V_{LN}$   & Line-to-neutral (phase-ground) voltage \\
  \end{tabular}
\end{frame}

% ---- Abbreviations ----
\begin{frame}{Abbreviations}
  \renewcommand{\arraystretch}{1.3}
  \begin{tabular}{@{}p{1.5cm}p{9.0cm}@{}}
    \textbf{Abbr.} & \textbf{Meaning} \\
    \hline
    AC   & Alternating current \\
    DC   & Direct current \\
    HV   & High Voltage \\
    HVDC & High-voltage direct current \\
    OHL  & Overhead line \\
  \end{tabular}
\end{frame}


\begin{frame}{Current calculation}
  \begin{itemize}
    \setlength{\itemsep}{0.5em}
    \item The current in the on the DC line or busbar is calculated by:
    \[ I = \frac{P}{\sqrt{3}\cdot V_{LL}} \]
    \item Note that for current calculation the phase-phase voltage $V_{LL}$ is used
    \item The maximum current is limited not only by the DC line capability but also by equipment on the DC line such as circuit breakers. 
    The maximum current rating for circuit breakers is $5000A$ (Hitachi Energy, GE).

  \end{itemize}
\end{frame}